Cette thèse, intitulée ``Rejeu et modélisation spatio-temporelle basée sur les
règles d'éléments archéologiques'', a pour objet la conception et le
développement d'un logiciel de modélisation 4D (3D+temps) permettant la
construction de bâtiments anciens pouvant évoluer au cours du temps,
l'enregistrement du processus de modélisation et le rejeu automatique de ce
processus afin de pouvoir générer plusieurs variantes du même bâtiment. %

In fine, ces restitutions virtuelles pourront constituer des documents de
travail pour des projets de restauration de vestiges. %

D'un point de vue technique, cela consiste en plusieurs objectifs. %
Enregistrer le processus de modélisation. %
Rejouer (ou réévaluer) de manière automatique des modèles. %
Faciliter la création de variantes de modèles. %
Visualiser l'évolution des modèles sur une échelle temporelle. %
Annoter des modèles à l'aide de données historiques. %

L'enregistrement du processus de modélisation dans une spécification
paramétrique est un pré-requis au rejeu automatique. %
Une spécification paramétrique est composée d'un ensemble d'opérations de
modélisation paramétrée appliquées successivement sur le modèle. %
Les paramètres de ces opérations sont classiquement des paramètres géométriques
ou topologiques, c'est-à-dire des références à des entités topologiques. %

%%% Local Variables:
%%% mode: latex
%%% TeX-master: "../main"
%%% End:
